\documentclass[10pt]{article}
\usepackage[T1]{fontenc}
\usepackage[utf8]{inputenc}
\usepackage[french]{babel}
\usepackage{amsmath,amssymb,mathrsfs,mathpazo}
\usepackage{geometry}
\geometry{a4paper,margin=2cm}
\usepackage{graphicx,array,tabularx,tabularray,colortbl,multirow,multicol,booktabs,enumitem,pifont,ulem,fancybox,diagbox,cancel}
\usepackage{pstricks,pstricks-add,pst-plot,pst-tree,pst-node,pst-text,pst-func,pst-math,pst-bspline,pst-3dplot,pst-eucl,pst-solides3d}
\let\pscalculate\undefined
\usepackage{pst-bezier}
\usepackage{tikz,pgfplots,tkz-tab}
\usetikzlibrary{arrows,arrows.meta,patterns,calc,positioning,trees,decorations.markings,shapes.geometric,shapes.misc,shapes.arrows,intersections,decorations.pathmorphing,decorations.pathreplacing,backgrounds,shadows,babel}
\pgfplotsset{compat=1.18}
\usepgfplotslibrary{fillbetween}
\usepackage[np]{numprint}
\usepackage{listings,makecell,pas-tableur,twoopt}
\let\barre\undefined
\usepackage[pstricks]{tablvar}
\definecolor{darkgreen}{RGB}{0,100,0}
\providecommand{\red}{\color{red}}
\providecommand{\blue}{\color{blue}}
\providecommand{\green}{\color{green}}
\providecommand{\black}{\color{black}}
\providecommand{\white}{\color{white}}
\providecommand{\gray}{\color{gray}}
\providecommand{\cyan}{\color{cyan}}
\providecommand{\magenta}{\color{magenta}}
\providecommand{\yellow}{\color{yellow}}
\providecommand{\violet}{\color{violet}}
\providecommand{\orange}{\color{orange}}
\providecommand{\euro}{\texteuro}
\providecommand{\EUR}{\texteuro}
\providecommand{\ieme}{\textsuperscript{e}}
\providecommand{\degres}{\ensuremath{^\circ}}
\providecommand{\degree}{\ensuremath{^\circ}}
\providecommand{\vv}[1]{\overrightarrow{#1}}
\providecommand{\Coord}[1]{#1}
\providecommand{\decosix}{}
\providecommand{\starredbullet}{\ensuremath{\bullet}}
\pagestyle{empty}
\setlength{\parindent}{0pt}
\usepackage{adjustbox}\DeclareUnicodeCharacter{2236}{:}\DeclareUnicodeCharacter{2190}{\ensuremath{\leftarrow}}
\providecommand{\R}{\mathbb{R}}
\providecommand{\N}{\mathbb{N}}
\providecommand{\Z}{\mathbb{Z}}
\providecommand{\Q}{\mathbb{Q}}
\providecommand{\D}{\mathbb{D}}
\providecommand{\C}{\mathbb{C}}
\providecommand{\e}{{\mathrm{e}}}
\providecommand{\vect}[1]{\overrightarrow{#1}}
\providecommand{\vectt}[1]{\overrightarrow{\mathrm{#1}}}
\providecommand{\vvt}[1]{\overrightarrow{\mathrm{#1}}}
\providecommand{\np}[2][]{{#2~#1}}
\providecommand{\vv}[1]{{\overrightarrow{#1}}}
\renewcommand{\d}{{\mathrm{d}}}
\providecommand{\degres}{{^\circ}}
\providecommand{\degre}{{^\circ}}
\providecommand{\blue}{\color{blue}}
\providecommand{\red}{\color{red}}
\providecommand{\enskip}{\quad}
\providecommand{\llbracket}{\lbrack\!\lbrack}
\providecommand{\rrbracket}{\rbrack\!\rbrack}
\providecommand{\euro}{{\text{€}}}
\providecommand{\no}{{\text{n°}}}
\providecommand{\interval}[2]{{[#1;#2]}}
\providecommand{\intervFO}[2]{{[#1;#2[}}
\providecommand{\dotfill}{\ldots}

\providecommand{\ds}{\displaystyle}
\providecommand{\textding}[1]{#1}
\providecommand{\Oij}{\(\mathrm{O};\vec{\imath},\vec{\jmath}\)}
\providecommand{\Oijk}{\(\mathrm{O};\vec{\imath},\vec{\jmath},\vec{k}\)}
\providecommand{\Ouv}{\(\mathrm{O};\vec{u},\vec{v}\)}
\providecommand{\intervFF}[2]{\ensuremath{[#1;#2]}}
\renewcommand{\euro}{\ifmmode\text{\texteuro}\else\texteuro\fi}
\DeclareUnicodeCharacter{00B0}{\ifmmode^{\circ}\else\textdegree\fi}
\DeclareUnicodeCharacter{20AC}{\ifmmode\text{\texteuro}\else\texteuro\fi}
% Énoncés : Éducation nationale. Transcription et corrigés : APMEP, Denis Vergès et contributeurs. Découpage et classement : CoopMaths, Sébastien Lozano, Éric Elter et contributeurs. Adaptation de la présentation : Maths Explorer.
% CC BY-SA : https://coopmaths.fr/www/about/
% https://coopmaths.fr/alea/static/e3c/2020/tex/e3c_2020_00_sujet20_2_cor.tex
\begin{document}
\noindent\textit{APMEP / CoopMaths. CC BY-SA. Présentation adaptée par Maths Explorer.}\par\medskip
\subsection*{1.}
	
	\subsubsection*{a.	$C(5) = 5^3 - 18 \times 5^2 + 750 \times 5 + 200 = 3625 $ euros.}
	
	\subsubsection*{b.}
	Le coût moyen de fabrication d’un millier d’objets lorsqu’on fabrique 5 000 objets est donc :
	\[
	\dfrac{3625}{5} = 725 \, \text{euros}.
	\]
	
	\subsection*{2.}
	
	Le coût moyen \(CM(q)\) de fabrication de \(q\) milliers d’objets, exprimé en euros, est donné par l’expression :
	
	\[
	CM(q) = \dfrac{C(q)}{q} = q^2 - 18q + 750 + \dfrac{200}{q}.
	\]
	
	\subsubsection*{a.}
		On a sur l’intervalle \([1 ; 20]\),
	\[
	C'_M(q) = 2q - 18 - \dfrac{200}{q^2} = \dfrac{2q^3 - 18q^2 - 200}{q^2}.
	\]
	
	Or
\begin{align*}
	2(q - 10)(q^2 + q + 10) &= (2q - 20)(q^2 + q + 10)\\
	& = 2q^3 + 2q^2 + 20q - 20q^2 - 20q - 200 \\
	&= 2q^3 - 18q^2 - 200
\end{align*}

	
	
	soit le numérateur de \(C'_M(x)\). \\
	Donc, pour tout \(q \in [1 ; 20]\),
	\[
	C'_M(q) = \dfrac{2(q - 10)(q^2 + q + 10)}{q^2}.
	\]
	
	\subsubsection*{b.}
	

	Pour le trinôme \(q^2 + q + 10\), on a \(\Delta = 1 - 40 = -39\) : pas de racines réelles et on sait que le trinôme est positif (signe de 1, coefficient de \(q^2\)).
	
	Comme \(q^2 > 0\) sur \([1 ; 20]\), le signe de \(C'_M(q)\) est celui de \(q - 10\).
	
	Donc sur \([1 ; 10]\), \(C'_M(q) < 0\) :\\
	 La fonction \(C_M\) est décroissante de \(C_M(1) = 1 - 18 + 750 + 200 = 933\)\\
	  à \(C_M(10) = 10^2 - 18 \times 10 + 750 + 20 = 690\).
	
	Sur \([10 ; 20]\), la fonction \(C_M\) est croissante de \(CM(10) = 690\) à \(CM(20) = 20^2 - 18 \times 20 + 750 + 10 = 800\).
	
	\begin{center}
	\begin{tikzpicture}[double distance=2pt]
		\tkzTabInit{$x$/1,$C'_M(x)$/1,$C_M$/3}{$-\infty$,$10$, $+\infty$}
		\tkzTabLine{,-,z,+}
		\tkzTabVar{+/$933$,-/$690$,+/$800$}
	\end{tikzpicture}
\end{center} 
	
	\subsubsection*{c.}
	
	Quel est le coût moyen minimal et pour quelle quantité d’objets est-il obtenu ?
	
	On voit sur le tableau que le coût moyen minimal de 690 est obtenu lorsque l’entreprise fabrique 10 000 objets.

\end{document}
